% Test of footnote reuse across the pages of a spread (didactic, twoside): a
% citation's footnote is set once per spread, and a later citation of the same
% work on the same spread has only the mark, with the number of the note.  See
% test-footreuse-spread.py, which checks the outcome from the PDF: every mark has
% a note with its number on its spread, and no spread has a note twice.  Built
% with \def\testcase{...} (and \def\testclass{memoir} for notes in the margin):
%   same    page 2 (verso) cites A first, page 3 (recto) again: one note, on 2,
%           and the mark on 3 has its number;
%   cross   page 3 (recto) cites B first, page 4 (verso) again: a new spread, so a
%           new number and a note on 4;
%   parted  page 2 cites A; on page 3 a later citation of A shares a paragraph
%           with a footnote too tall for the rest of the page, so its line goes to
%           page 4, another spread: it ends with a new number and its own note
%           there (a repeat of the old number exists only in the run that finds
%           the parting).
% With memoir, a note counts as on the page where it starts (the margin of page 3
% is full, so the notes go to its foot):
%   splitreuse  page 3 cites A twice on one line near its foot; the note is split
%           over pages 3 and 4 (spreads 2--3 and 4--5): the second mark reuses
%           it, since it starts on page 3;
%   splitrest   page 3 cites A, its note split over pages 3 and 4; page 4 cites A
%           again, beside the rest: a new spread, so a new number and a note;
%   overreuse   page 3 cites B, whose note runs past the foot of the margin, and
%           again lower on the page: the second mark reuses it;
%   waitreuse   on page 3 a long footnote is split at the foot, and B is cited
%           twice on its line: B's note waits for the foot of page 4, and the
%           second mark reuses it (both marks on page 3, the note on 4, in order);
%   carryreuse  page 3 cites B twice on its last line, where the foot has no
%           room: the note is carried to page 4, and the second mark reuses it.
% (In these two the margin's tall note runs past its foot itself, so that B's
% note can't: it goes to the foot.)
%   chapreuse  memoir restarts footnote numbers in each chapter: chapter One
%           cites B twice on a page, the second mark reusing note 1; in chapter
%           Two, note 1 (a plain footnote, no citation) is on its page's last
%           line, where the foot has no room, and is carried: its warning names
%           no mark, since nothing reuses chapter Two's note 1.
\ifdefined\testclass\else\def\testclass{article}\fi
\ifdefined\testcase\else\def\testcase{same}\fi
\documentclass[twoside]{\testclass}
\usepackage[backend=biber,autocite=footnote,citestyle=verbose,style=verbose]{biblatex}
\usepackage{csquotes}
% \testfloor: the same document with noteoverflow=carry, for the runs LaTeX needs
% without moving any note (the floor of the runs to settle).
\ifdefined\testfloor\PassOptionsToPackage{noteoverflow=carry}{didactic}\fi
\usepackage{didactic}
\addbibresource{test-footcite.bib}
\def\same{same}\def\cross{cross}\def\parted{parted}
\def\splitreuse{splitreuse}\def\splitrest{splitrest}\def\overreuse{overreuse}
\def\waitreuse{waitreuse}\def\carryreuse{carryreuse}\def\chapreuse{chapreuse}
\newcommand\longpost{with a long postnote that makes the note run over several
  lines of the foot, as a note with a quotation in it may well do, and then some
  more words, so that it can't fit the room left at the foot of its page, nor
  the room and a line, and has to be split over two pages: \filler\ \filler\
  \filler\ \filler}
\newcommand\filler{Some words to fill up the rest of this line of the text and
  also quite a few of the next one, as the words of a paragraph do.}
\begin{document}
\ifx\testcase\splitreuse\def\testsplit{}\fi
\ifx\testcase\splitrest\def\testsplit{}\fi
Page~1 (recto), intentionally left without citations.
\newpage
\ifx\testcase\cross\def\testnotwo{}\fi
\ifx\testcase\chapreuse\def\testnotwo{}\fi
\ifdefined\testnotwo
  Page~2 (verso), without citations.
\else
  Page~2 (verso): the first citation of the spread\footcite{A}.
\fi
\newpage
\ifx\testcase\same
  Page~3 (recto): the same spread as page~2, so this\footcite{A} has the number
  of the note on page~2.
\fi
\ifx\testcase\cross
  Page~3 (recto): the first citation of B\footcite{B}.
  \newpage
  Page~4 (verso): a new spread, so this\footcite{B} has a new number.
\fi
\ifdefined\testsplit
  \noindent x\marginpar{\rule{1pt}{.99\textheight}}%
  \par\noindent\rule{1cm}{\dimexpr\textheight-6\baselineskip\relax}\par
  Page~3, near its foot: A\footcite[\longpost]{A}%
  \ifx\testcase\splitreuse\ and A again\footcite[\longpost]{A}\fi
  \ some more words here.
  \ifx\testcase\splitrest
    \newpage
    Page~4: A again\footcite[\longpost]{A}, beside the rest of its note.
  \fi
\fi
\ifx\testcase\overreuse
  \noindent x\marginpar{\rule{1pt}{.98\textheight}}%
  \par\noindent\rule{1cm}{\dimexpr\textheight-12\baselineskip\relax}\par
  Page~3, near its foot: B\footcite{B}, and then again B\footcite{B} on the
  same page.
\fi
\ifx\testcase\waitreuse
  \noindent x\marginpar{\rule{1pt}{1.01\textheight}}%
  \par\noindent\rule{1cm}{\dimexpr\textheight-6\baselineskip\relax}\par
  Page~3, near its foot: a long note\footnote{LONGNOTE \longpost}, then
  B\footcite{B} and B again\footcite{B}, all on this line.
\fi
\ifx\testcase\carryreuse
  \noindent x\marginpar{\rule{1pt}{1.01\textheight}}%
  \par\noindent\rule{1cm}{\dimexpr\textheight-2\baselineskip\relax}\par
  Page~3, its last line: B\footcite{B} and B again\footcite{B}.
\fi
\ifx\testcase\chapreuse
  \chapter{One}
  The first citation of B\footcite{B}, and B again\footcite{B} on the same
  page.
  \chapter{Two}
  A page of its own.
  \newpage
  \noindent x\marginpar{\rule{1pt}{1.01\textheight}}%
  \par\noindent\rule{1cm}{\dimexpr\textheight-2\baselineskip\relax}\par
  Its last line: a note\footnote{CHAPNOTE, carried to the next page.}.
\fi
\ifx\testcase\parted
  \noindent\rule{1cm}{\dimexpr\textheight-4\baselineskip\relax}\par
  Words with a tall footnote\footnote{TALLNOTE\rule{0pt}{8\baselineskip}}
  \filler\ And then the later citation of A\footcite{A}, which goes to page~4
  with its line.
\fi
\end{document}
